\documentclass[a4paper,12pt,oneside]{book}
\linespread{1.5}
\usepackage{geometry}
\geometry{top=3cm,bottom=3cm,left=3.5cm,right=3.5cm} % Impostiamo il valore dei margini.
\usepackage[utf8]{inputenc}
\usepackage{tikz,pgf}
\usepackage{indentfirst}
\usepackage{amsfonts}
\usepackage[italian]{babel} %% Abbiamo istruito LaTeX ad usare la lingua italiana.
\usepackage{todonotes}
\usepackage{subfig}
\usepackage{fancyhdr}
\setlength{\parindent}{1cm}
\usepackage{graphicx}
\usepackage{float}
\usepackage{csquotes}
\usepackage{amsmath}
\usepackage{numprint}
\usepackage{pgfplots}
\usepackage{pgfplotstable}
\usepackage{tocbibind}

\graphicspath{{imgs/}} %% cartella dalla quale ricavare le immagini %%

%%%%%%%%%%%%%%%%%%%%%%%%%%%%%%%%%%%%%%%%%%%%%%%%%%%%%%%%%%%%
%% Template per la tesi di laurea triennale dell' UniCal  %%
%% Autore: Fabio gimondo                                  %%
%% Data: Aprile 2020                                      %%
%%%%%%%%%%%%%%%%%%%%%%%%%%%%%%%%%%%%%%%%%%%%%%%%%%%%%%%%%%%%

\setlength {\marginparwidth }{2cm} 
\pgfplotsset{compat=1.16}

\newenvironment{dedication}
  {\clearpage           % we want a new page
   \thispagestyle{empty}% no header and footer
   \vspace*{\stretch{1}}% some space at the top 
   \itshape             % the text is in italics
   \raggedleft          % flush to the right margin
  }
   {\par % end the paragraph
   \vspace{\stretch{3}} % space at bottom is three times that at the top
   \clearpage           % finish off the page
  }
  
\begin{document}
%% Frontespizio %%
\begin{titlepage}
    \begin{center}
        \LARGE{\uppercase{\textbf{Università della Calabria}}}\\
        \vspace{5mm}
%% Dipartimento %%
    	\uppercase{\normalsize Dipartimento di Ingegneria del LaTeX}\\
    \end{center}
    \begin{figure}[H]
        \centering
        \includegraphics[width=0.7\textwidth]{imgs/logo.png}
    \end{figure}
    
\begin{tikzpicture}[remember picture,overlay]
\node [opacity=0.1,scale=1.15] at (current page.center) {\includegraphics{logo1.png}};
\end{tikzpicture}
    
    \begin{center}
%% Corso di Laurea %%
    	\large{ Corso di Laurea in \textit{LaTeX}ologia}\\
    	\vspace{7mm}
%% Titolo della Tesi %%
        {\LARGE{\bf Titolo Tesi}}\\
    	\vspace{3mm}
    \end{center}
    
    \vspace{15mm}
    \noindent
    \begin{minipage}[t]{0.47\textwidth}
%% Relatore %%
    	{\large{ Relatore:\\\bf prof. \\Nome Cognome}}
    	\vspace{12mm}\\
%% Correlatore (qualora non ci fosse, anteporre "%" ai codici relativi) 
    	{\large{ Correlatore:\\ \bf prof. \\Nome Cognome}}
    \end{minipage}
    \hfill
    \begin{minipage}[t]{0.4\textwidth}\raggedleft
%% Candidato %%
    	{\large{Candidato: \\ \bf Nome Cognome\\ Mat. 123456}}
    \end{minipage}
    
    \vspace{20mm}
    
%% Anno Accademico %%
    \centering{\large \uppercase{ Anno Accademico 20xx/20yy}}

\end{titlepage}

\begin{dedication}
Dedicata a coloro che\\ ogni giorno lottano per i propri sogni.
\end{dedication}

%% Corpo della tesi %%

\chapter*{Sommario} %% In questa zona inseriamo l'abstract del nostro documento. %%
Inserisci abstract: ...


\tableofcontents
\input{capitoli/capitolo1}
\input{capitoli/capitolo2}
\input{capitoli/capitolo3}
\input{capitoli/capitolo4}
\input{capitoli/capitolo5}
\input{capitoli/capitolo6}
%Eliminare/Aggiungere/Personalizzare i capitolo nella cartella "capitoli"

%\chapter*{Conclusioni} In questo modo non comparirà dell'indice

\begin{thebibliography}{100} %% 100 rappresenta il n° massimo di voci all'interno della bibliografia: può essere modificato a vostro piacimento. %%
\bibitem{[1]:}
\bibitem{[2]:}
\bibitem{[3]:}
\bibitem{[4]:}

\end{thebibliography}

\chapter*{Ringraziamenti}
Inserisci frasi sdolcinate: ...
\end{document}
